% Apuntes de teoría en PDF · 4.º ESO · Unidad 9: Estadística (versión en castellano)
% Mateix estil que els apunts de Batxillerat (franja i caixes de teoria en blau, exemples en roig), però més visuals,
% curts i amb dibuixos, per a l'alumnat de 3r ESO. Es compila amb pdflatex dues vegades (per als enllaços de l'índex),
% amb logo-fentmates.pdf (_eines/logo/pdf/horitzontal-clar.pdf) al costat.
\documentclass[11pt,a4paper]{article}
\usepackage[utf8]{inputenc}
\usepackage[T1]{fontenc}
\usepackage[spanish,es-noshorthands,es-nodecimaldot,es-nolists]{babel}
\usepackage[sfdefault,book,lining]{FiraSans}
\usepackage[italic,noplusnominus,noequal]{mathastext}
\linespread{1.18}
\usepackage[a4paper,top=2.4cm,bottom=2.2cm,left=1.6cm,right=1.6cm,headheight=1.1cm,headsep=0.5cm,footskip=1cm]{geometry}
\usepackage{amsmath,amssymb,cancel}
\usepackage{icomma}
\usepackage{graphicx}
\usepackage[table]{xcolor}
\usepackage{tikz}\usetikzlibrary{calc,positioning,arrows.meta,decorations.pathreplacing,shapes.geometric}
\usepackage{fontawesome5}
\usepackage{ccicons}
\usepackage{fancyhdr}
\usepackage{enumitem}
\usepackage{array}
\usepackage{needspace}
\usepackage[most]{tcolorbox}
\usepackage{titlesec}
\usepackage[hidelinks]{hyperref}
\hypersetup{pdftitle={Apuntes de teoría. Unidad 9: Estadística. 4.º ESO},pdfsubject={Matemáticas A 4.º ESO},pdfkeywords={estadística, población, muestra, variables estadísticas, frecuencias, tablas, intervalos, marca de clase, gráficos estadísticos, media, mediana, moda, rango, desviación típica, coeficiente de variación, cuartiles, percentiles, diagrama de cajas}}

% colors de la web (bloc d'Àlgebra: blau)
\definecolor{fmnavy}{HTML}{1B2748}
\definecolor{fmblue}{HTML}{3D6AA0}
\definecolor{fmbluel}{HTML}{4E81BD}
\definecolor{fmyellow}{HTML}{FFC94D}
\definecolor{fmgold}{HTML}{E9A919}
\definecolor{fmaccent}{HTML}{FF7A50}
\definecolor{fmdef}{HTML}{1D4F73}\definecolor{fmdefbg}{HTML}{EEF4F9}
\definecolor{fmex}{HTML}{A3203A}\definecolor{fmexbg}{HTML}{FBF0F2}
\definecolor{fmsf}{HTML}{C2410C}
\definecolor{fmok}{HTML}{15803D}
\definecolor{fmmuted}{HTML}{5D6685}
\definecolor{fmsurface}{HTML}{EEF1F8}
\definecolor{fmA}{HTML}{2F6FD6}\definecolor{fmB}{HTML}{EA6A10}\definecolor{fmR}{HTML}{E11D48}\definecolor{fmK}{HTML}{16A34A}\definecolor{fmV}{HTML}{7C3AED}

% lletres
\newcommand{\comic}{\fontfamily{ComicNeue-TLF}\selectfont}
\newcommand{\heavy}{\fontseries{eb}\selectfont}
\newcommand{\semi}{\fontseries{sb}\selectfont}
\DeclareUnicodeCharacter{2212}{-}

\setlength{\parindent}{0pt}
\setlength{\parskip}{7pt}
\setlength{\lineskip}{4pt}\setlength{\lineskiplimit}{4pt}
\setlength{\emergencystretch}{3em}
\setlist{itemsep=3pt,topsep=3pt,leftmargin=1.4em}
\setcounter{secnumdepth}{1}
\renewcommand{\thesection}{9.\arabic{section}}
\titleformat{\section}[hang]{\color{fmnavy}\fontsize{23}{27}\heavy\selectfont}{\tcbox[on line,colback=fmblue,colframe=fmblue,arc=3mm,boxsep=2pt,left=6pt,right=6pt,top=3pt,bottom=3pt]{\color{white}\thesection}}{0.45em}{}[\vspace{5pt}{\color{fmblue}\titlerule[2pt]}]
\titlespacing*{\section}{0pt}{4pt}{6pt}
\titleformat{\subsection}{\color{fmblue}\Large\semi}{}{0pt}{\faCaretRight\ }
\titlespacing*{\subsection}{0pt}{8pt}{2pt}

% frase de cada apartat, «a mà»
\newcommand{\frase}[1]{\par\vspace{-2pt}{\comic\Large\color{fmbluel}\faCommentDots[regular]\ \ #1\par}\vspace{2pt}}
\newcommand{\seccio}[3]{\Needspace{16\baselineskip}\section{#2}\hypertarget{#1}{}\frase{#3}}
% notes a mà dins dels exemples
\newcommand{\nota}[1]{{\comic\small\color{fmex!85!black}#1}}
\newcommand{\notab}[1]{{\comic\small\color{fmblue}#1}}
% bé / malament
\newcommand{\be}[1]{{\color{fmok}\faCheckCircle}\ \ensuremath{#1}}
\newcommand{\mal}[1]{{\color{fmR}\faTimesCircle}\ {\color{fmR}\ensuremath{#1}}}
% pas numerat
\newcommand{\pas}[1]{\tikz[baseline=(p.base)]\node[circle,fill=fmblue,text=white,inner sep=1.2pt,minimum size=1.25em,font=\small\heavy](p){#1};}
% etiqueta de color
\newcommand{\chip}[2]{\tcbox[on line,colback=#1!14,colframe=#1!60,arc=2mm,boxsep=1pt,left=5pt,right=5pt,top=1pt,bottom=1pt]{#2}}
\newcommand{\fmkey}[1]{\colorbox{fmyellow!55}{\strut #1}}
\newcommand{\cA}[1]{{\color{fmA}#1}}\newcommand{\cB}[1]{{\color{fmB}#1}}\newcommand{\cK}[1]{{\color{fmK}#1}}\newcommand{\cR}[1]{{\color{fmR}#1}}\newcommand{\cV}[1]{{\color{fmV}#1}}

% caixes
\tcbset{fmbase/.style={enhanced,arc=3mm,boxrule=1.2pt,left=4mm,right=4mm,top=2mm,bottom=2.5mm,before skip=8pt,after skip=8pt,fonttitle=\heavy\large,before upper=\raggedright}}
\newtcolorbox{idea}[1][Idea clave]{fmbase,colback=fmdefbg,colframe=fmdef,coltitle=white,title={\faKey\ \ #1}}
\newtcolorbox{exemple}[1][Ejemplo]{fmbase,colback=fmexbg,colframe=fmex,colbacktitle=fmexbg,coltitle=fmex,title={\faPencil*\ \ #1},titlerule=0.5pt,titlerule style={fmex!40,dashed}}
\newtcolorbox{compte}[1][¡Cuidado!]{fmbase,colback=fmsf!6,colframe=fmsf,coltitle=white,fonttitle=\comic\bfseries\Large,title={\faExclamationTriangle\ \ #1}}
\newtcolorbox{truc}[1][Truco]{fmbase,colback=fmyellow!20,colframe=fmgold,colbacktitle=fmyellow,coltitle=fmnavy,fonttitle=\comic\bfseries\Large,title={\faLightbulb[regular]\ \ #1}}

% dibuix de la targeta de la unitat, en blanc
\newcommand{\fmcard}{\begin{tikzpicture}[x=0.062cm,y=-0.062cm,white,line width=3pt,line cap=round]
  \begin{scope}[transparency group,opacity=0.15]
  \node[anchor=base west,inner sep=0pt,font=\fontsize{28}{28}\selectfont] at (14,44) {$\bar x\quad\sigma$};
  \draw[line width=2pt] (120,110) -- (235,110);
  \fill (130,110) rectangle (148,70); \fill (155,110) rectangle (173,30); \fill (180,110) rectangle (198,55); \fill (205,110) rectangle (223,85);
  \end{scope}
\end{tikzpicture}}
\newcommand{\fmbrand}{\raisebox{-0.25\height}{\includegraphics[height=0.5cm]{logo-fentmates.pdf}}}

% regla de Ruffini: #1 = a, #2 = coeficients, #3 = fila del mig (sense el primer), #4 = fila de baix, #5 = nombre de coeficients
\newcommand{\ruffini}[5]{\begin{tikzpicture}[x=1.25cm,y=0.85cm,every node/.style={font=\large}]
  \foreach \v [count=\i] in {#2} \node (t\i) at (\i,0) {$\v$};
  \foreach \v [count=\i] in {#3} {\pgfmathtruncatemacro\j{\i+1}\node (m\j) at (\j,-1) {$\v$};}
  \foreach \v [count=\i] in {#4} \node (b\i) at (\i,-2.1) {$\v$};
  \node[text=fmB] at (0,-1) {$#1$};
  \draw[line width=1.1pt,fmnavy] (0.45,0.5)--(0.45,-1.5) (0.45,-1.5)--(#5+0.5,-1.5);
  \node[draw=fmR,line width=1.3pt,rounded corners=2mm,minimum size=0.75cm] at (#5,-2.1) {};
  \draw[-{Latex[length=2.2mm]},fmK,line width=1.1pt] (t1.south) -- (b1.north);
  \pgfmathtruncatemacro\last{#5-1}
  \foreach \i in {1,...,\last} {\pgfmathtruncatemacro\j{\i+1}
    \draw[-{Latex[length=2.2mm]},fmB,line width=1.1pt,shorten >=1pt] (b\i.north east) -- (m\j.south west);
    \draw[-{Latex[length=2.2mm]},fmA!70,line width=1pt,dashed,shorten >=2pt] ([xshift=8pt]t\j.south) -- ([xshift=8pt]b\j.north);}
\end{tikzpicture}}

% capçalera i peu
\pagestyle{fancy}\fancyhf{}
\fancyhead[L]{\small\fmbrand}
\fancyhead[R]{\small\color{fmmuted}4.º ESO · Unidad 9: Estadística}
\fancyfoot[L]{\scriptsize\color{fmmuted}\href{https://creativecommons.org/licenses/by-nc-sa/4.0/deed.es}{\ccbyncsa}\ \ FentMates · fentmates.org · CC BY-NC-SA 4.0}
\fancyfoot[R]{\small\color{fmmuted}\thepage}
\renewcommand{\headrulewidth}{0.4pt}\renewcommand{\headrule}{\color{fmmuted!40}\hrule height \headrulewidth}
\fancypagestyle{primera}{\fancyhf{}\renewcommand{\headrulewidth}{0pt}
  \fancyfoot[L]{\scriptsize\color{fmmuted}\href{https://creativecommons.org/licenses/by-nc-sa/4.0/deed.es}{\ccbyncsa}\ \ FentMates · fentmates.org · CC BY-NC-SA 4.0}
  \fancyfoot[R]{\small\color{fmmuted}\thepage}}

% targeta de l'índex
\newcommand{\itile}[4]{\begin{tcolorbox}[enhanced,colback=white,colframe=fmblue!35,arc=3mm,boxrule=1pt,boxsep=0pt,left=6pt,right=6pt,top=5pt,bottom=5pt,before skip=0pt,after skip=7pt]
  \hyperlink{#1}{\begin{minipage}[c]{0.9cm}\centering{\color{fmblue}\LARGE #3}\end{minipage}\hspace{6pt}%
  \begin{minipage}[c]{0.78\linewidth}{\color{fmblue}\heavy #2}\\[-2pt]{\color{fmnavy}\semi\large #4}\end{minipage}}\end{tcolorbox}}

% taula amb vores arrodonides: el contingut es retalla amb la forma (sense @{} a les vores ni entre columnes)
\newsavebox\tabclipbox
\newcommand{\taulaclip}[3][3mm]{\sbox\tabclipbox{#3}\begin{tikzpicture}[baseline=(tcn.base)]
  \begin{scope}\clip[rounded corners=#1] (0,-\dp\tabclipbox) rectangle (\wd\tabclipbox,\ht\tabclipbox);
  \node[inner sep=0pt,outer sep=0pt,anchor=base west] (tcn) at (0,0) {\usebox\tabclipbox};\end{scope}
  \draw[#2,rounded corners=#1] (0,-\dp\tabclipbox) rectangle (\wd\tabclipbox,\ht\tabclipbox);
\end{tikzpicture}}
\newcommand{\taularod}[3][fmblue]{\taulaclip{#1,line width=1.2pt}{\renewcommand{\arraystretch}{1.5}\begin{tabular}{#2}#3\end{tabular}}}
% balança
\newcommand{\balanca}[3]{\begin{tikzpicture}[scale=#3,transform shape,every node/.style={anchor=center}]
  \fill[fmnavy] (-0.45,-1.9) rectangle (0.45,-1.75);
  \fill[fmnavy!80] (-0.12,-1.75) rectangle (0.12,0);
  \fill[fmgold] (0,0.05) circle (0.13);
  \draw[fmnavy,line width=2.4pt,line cap=round] (-2,0) -- (2,0);
  \foreach \s in {-1,1} {
    \draw[fmmuted,line width=0.8pt] (\s*2,0) -- (\s*2-0.75,-1.05) (\s*2,0) -- (\s*2+0.75,-1.05);
    \fill[fmyellow!70,draw=fmgold,line width=1pt] (\s*2-0.95,-1.05) arc[start angle=180,end angle=360,x radius=0.95,y radius=0.32] -- cycle;}
  \node[font=\Large,fill=fmA!15,draw=fmA,rounded corners=1.5mm,inner sep=3pt] at (-2,-0.68) {$#1$};
  \node[font=\Large,fill=fmB!15,draw=fmB,rounded corners=1.5mm,inner sep=3pt] at (2,-0.68) {$#2$};
\end{tikzpicture}}
% eixos amb quadrícula (dins d'un tikzpicture): xmin, xmax, ymin, ymax
\newcommand{\eixos}[4]{\draw[fmmuted!22,line width=0.4pt] (#1,#3) grid (#2,#4);
  \draw[-{Latex[length=2mm]},fmmuted,line width=0.8pt] (#1-0.2,0)--(#2+0.4,0) node[right,font=\small]{$x$};
  \draw[-{Latex[length=2mm]},fmmuted,line width=0.8pt] (0,#3-0.2)--(0,#4+0.4) node[above,font=\small]{$y$};}
% taula de traduccions (diccionari) amb les vores arrodonides
\newcommand{\dic}[2]{#1 & {\large$#2$}\\}
\newcommand{\dicratio}{0.62}
\newcommand{\dichead}{\begin{tabular}{@{}>{\raggedright}p{\dimexpr\dicratio\linewidth-12pt}>{\centering\arraybackslash}p{\dimexpr\linewidth-\dicratio\linewidth-12pt}@{}}\faComment[regular]\ Lo decimos así\dots & \dots\ así\end{tabular}}
\newenvironment{dicbox}{\begin{tcolorbox}[enhanced,arc=3mm,boxrule=1.2pt,colframe=fmblue,colback=white,coltitle=white,fonttitle=\heavy,
  title={\dichead},left=0pt,right=0pt,top=0pt,bottom=1.5mm,boxsep=0pt,lefttitle=6pt,righttitle=6pt,before skip=0pt,after skip=0pt]\rowcolors{1}{white}{fmsurface}%
  \begin{tabular}{@{\hspace{6pt}}>{\raggedright}p{\dimexpr\dicratio\linewidth-12pt}>{\centering\arraybackslash}p{\dimexpr\linewidth-\dicratio\linewidth-12pt}@{\hspace{6pt}}}}
  {\end{tabular}\end{tcolorbox}}
% resum
\newcommand{\rtile}[3]{\begin{tcolorbox}[enhanced,colback=white,colframe=fmblue!45,arc=3mm,boxrule=1pt,left=3mm,right=3mm,top=2.5mm,bottom=2.5mm,fontupper=\large,
  title={\tcbox[on line,colback=fmblue,colframe=fmblue,arc=2mm,boxsep=0pt,left=4pt,right=4pt,top=1pt,bottom=1pt]{\color{white}\small\heavy #1}\ \ #2},coltitle=fmnavy,colbacktitle=fmdefbg,fonttitle=\semi,toptitle=0.5mm,bottomtitle=0.5mm]#3\end{tcolorbox}}
% portada: #1 subtítol a mà
\newcommand{\portada}[1]{\thispagestyle{primera}
\begin{tikzpicture}[remember picture,overlay]
  \fill[fmblue] (current page.north west) rectangle ([yshift=-9.4cm]current page.north east);
  \begin{scope}
    \clip (current page.north west) rectangle ([yshift=-9.4cm]current page.north east);
    \node[anchor=north west,inner sep=0pt] at ([xshift=6.6cm,yshift=-1.3cm]current page.north west) {\fmcard};
  \end{scope}
  \node[anchor=north west,inner sep=0pt] at ([xshift=1.6cm,yshift=-1.1cm]current page.north west) {\tcbox[colback=white,colframe=white,arc=2.5mm,boxsep=0pt,left=7pt,right=7pt,top=5pt,bottom=5pt]{\includegraphics[height=0.75cm]{logo-fentmates.pdf}}};
\end{tikzpicture}
\vspace*{0.55cm}
\tcbox[colback=white!15!fmblue,colframe=white!15!fmblue,arc=2.5mm,boxsep=0pt,left=6pt,right=6pt,top=3pt,bottom=3pt]{\color{white}\small\bfseries Apuntes de teoría · 4.º ESO · Matemáticas A}

\vspace{2pt}
{\raggedright\color{white}\fontsize{30}{36}\heavy\selectfont{\color{fmyellow}Unidad 9}\\[2pt]Estadística\par}
\vspace{8pt}
{\color{white}\comic\Large #1\par}
\vspace*{\dimexpr 7.4cm-\pagetotal\relax}
}
\newcommand{\llegenda}{\par\vspace{16pt}
{\color{fmnavy}\heavy\Large\faPalette\ \ Cómo se leen estos apuntes}\par\smallskip
\begin{tcbraster}[raster columns=4,raster equal height,raster column skip=8pt,raster before skip=12pt]
\tcbox[colback=fmdefbg,colframe=fmdef,arc=3mm,boxrule=1.2pt,halign=center,width=\linewidth]{{\color{fmdef}\Large\faKey}\\[2pt]{\color{fmdef}\heavy Idea clave}\\[-1pt]{\small lo esencial}}
\tcbox[colback=fmexbg,colframe=fmex,arc=3mm,boxrule=1.2pt,halign=center,width=\linewidth]{{\color{fmex}\Large\faPencil*}\\[2pt]{\color{fmex}\heavy Ejemplo}\\[-1pt]{\small paso a paso}}
\tcbox[colback=fmsf!6,colframe=fmsf,arc=3mm,boxrule=1.2pt,halign=center,width=\linewidth]{{\color{fmsf}\Large\faExclamationTriangle}\\[2pt]{\color{fmsf}\heavy ¡Cuidado!}\\[-1pt]{\small errores típicos}}
\tcbox[colback=fmyellow!20,colframe=fmgold,arc=3mm,boxrule=1.2pt,halign=center,width=\linewidth]{{\color{fmgold!80!black}\Large\faLightbulb[regular]}\\[2pt]{\color{fmnavy}\heavy Truco}\\[-1pt]{\small para recordarlo}}
\end{tcbraster}}
\newcommand{\resumtitol}{\clearpage
\hypertarget{resum}{}
{\color{fmnavy}\fontsize{23}{27}\heavy\selectfont\faClipboardList\ \ En resumen}\par\vspace{5pt}{\color{fmblue}\titlerule[2pt]}
\frase{Toda la unidad en una página.}}
% els 5 passos per a resoldre problemes: #2 = text del pas 2
\newcommand{\passosproblema}[2]{\begin{center}
\begin{tikzpicture}[s/.style={draw=fmblue,fill=fmdefbg,rounded corners=3mm,line width=1.2pt,minimum width=2.75cm,minimum height=2.45cm,align=center,font=\small},
  ic/.style={font=\LARGE,text=fmblue}]
  \foreach \i/\ic/\t in {1/\faBookOpen/{\textbf{Leemos}\\y entendemos},2/\faQuestionCircle/{#1},3/\faEquals/{#2},4/\faCogs/{\textbf{Resolvemos}},5/\faCheckDouble/{\textbf{Comprobamos}\\y respondemos}} {
    \node[s] (s\i) at (3.3*\i,0) {};
    \node[ic] at ([yshift=-0.5cm]s\i.north) {\ic};
    \node[align=center,font=\small,anchor=north] at ([yshift=-0.95cm]s\i.north) {\t};
    \node[circle,fill=fmblue,text=white,font=\small\heavy,inner sep=1.5pt,minimum size=1.3em] at (s\i.north west) {\i};}
  \foreach \i [evaluate=\i as \j using int(\i+1)] in {1,...,4} \draw[-{Latex[length=3mm]},fmmuted,line width=1.5pt] (s\i) -- (s\j);
\end{tikzpicture}
\end{center}}

\newcommand{\lo}{\mathopen{]}}\newcommand{\ro}{\mathclose{[}}
% punts ple i buit, nombres dels eixos i taula de valors
\newcommand{\pple}[2][fmblue]{\fill[#1] #2 circle (3.6pt);}
\newcommand{\nums}[4]{\foreach \k in {#1,...,#2} {\ifnum\k=0\else\node[below,font=\scriptsize,text=fmmuted] at (\k,0) {$\k$};\fi}
  \foreach \k in {#3,...,#4} {\ifnum\k=0\else\node[left,font=\scriptsize,text=fmmuted] at (0,\k) {$\k$};\fi}}
\newcommand{\tv}[3][fmblue]{\taulaclip{#1,line width=1.2pt}{\renewcommand{\arraystretch}{1.45}%
  \begin{tabular}{>{\columncolor{#1}\color{white}\heavy}c*{12}{w{c}{0.8cm}}}#2\\\rowcolor{#1!8}\cellcolor{#1}#3\end{tabular}}}
% etiqueta d'una recta o corba
\newcommand{\etq}[3]{\node[fill=white,inner sep=1.5pt,text=#1,font=\small] at #2 {#3};}

% interval sobre la recta: #1 esquerre, #2 dret, #3 tipus esquerre (o/c/a), #4 tipus dret (o/c/a), #5 color
\def\intvdot#1#2#3{\ifx#2o\filldraw[fill=white,draw=#3,line width=1.3pt] (#1,0) circle (3.3pt);\fi\ifx#2c\fill[#3] (#1,0) circle (3.5pt);\fi}
\newcommand{\intv}[5]{\begin{tikzpicture}[x=0.4cm,baseline=-3pt]
  \draw[fmmuted,line width=0.8pt] (-6,0)--(7,0);
  \foreach \k in {-5,...,6} \draw[fmmuted,line width=0.6pt] (\k,-0.08)--(\k,0.08);
  \ifx#3a\def\xa{-6}\else\def\xa{#1}\fi
  \ifx#4a\def\xb{7}\else\def\xb{#2}\fi
  \ifx#4a\draw[#5,line width=2.6pt,-{Latex[length=2.6mm]}] (\xa,0)--(\xb,0);\else\ifx#3a\draw[#5,line width=2.6pt,{Latex[length=2.6mm]}-] (\xa,0)--(\xb,0);\else\draw[#5,line width=2.6pt] (\xa,0)--(\xb,0);\fi\fi
  \intvdot{#1}{#3}{#5}\intvdot{#2}{#4}{#5}
  \ifx#3a\else\node[below=2pt,font=\scriptsize] at (#1,0) {$#1$};\fi
  \ifx#4a\else\node[below=2pt,font=\scriptsize] at (#2,0) {$#2$};\fi
\end{tikzpicture}}

\newcommand{\passos}[1]{\begin{tcolorbox}[fmbase,colback=fmdefbg,colframe=fmdef,coltitle=white,title={\faListOl\ \ Cómo se resuelven}]#1\end{tcolorbox}}
\newcommand{\passosT}[2]{\begin{tcolorbox}[fmbase,colback=fmdefbg,colframe=fmdef,coltitle=white,title={\faListOl\ \ #1}]#2\end{tcolorbox}}
% barra d'un diagrama: #1 x central, #2 altura, #3 color
\newcommand{\barra}[3]{\filldraw[fill=#3!35,draw=#3,line width=0.9pt] (#1-0.32,0) rectangle (#1+0.32,#2);}
% diagrama de caixes: #1 y, #2 mín, #3 Q1, #4 Me, #5 Q3, #6 màx, #7 color
\newcommand{\caixa}[7]{\draw[#7,line width=1.3pt] (#2,#1) -- (#3,#1) (#5,#1) -- (#6,#1) (#2,#1-0.25) -- (#2,#1+0.25) (#6,#1-0.25) -- (#6,#1+0.25);
  \filldraw[fill=#7!20,draw=#7,line width=1.3pt] (#3,#1-0.45) rectangle (#5,#1+0.45); \draw[#7,line width=2.2pt] (#4,#1-0.45) -- (#4,#1+0.45);}


\begin{document}
\portada{Recogemos datos, los ordenamos, los dibujamos y los resumimos con unos pocos números.}

{\color{fmnavy}\heavy\Large\faMapSigns\ \ ¿Qué aprenderás?}\par\smallskip
\begin{minipage}[t]{0.485\linewidth}
\itile{sec1}{9.1}{\faUsers}{Población, muestra y variables}
\itile{sec2}{9.2}{\faTable}{Tablas de frecuencias}
\itile{sec3}{9.3}{\faLayerGroup}{Datos agrupados en intervalos}
\itile{sec4}{9.4}{\faChartBar}{Gráficos estadísticos}
\itile{sec5}{9.5}{\faBullseye}{Parámetros de centralización}
\end{minipage}\hfill
\begin{minipage}[t]{0.485\linewidth}
\itile{sec6}{9.6}{\faArrowsAltH}{Parámetros de dispersión}
\itile{sec7}{9.7}{\faMapMarker*}{Cuartiles y percentiles}
\itile{sec8}{9.8}{\faIcon{box}}{Diagrama de cajas}
\itile{resum}{\phantom{9.8}}{\faClipboardList}{En resumen}
\end{minipage}
\llegenda

% ============================================================ 9.1
\clearpage
\seccio{sec1}{Población, muestra y variables estadísticas}{Antes de contar, hay que saber a quién y qué preguntamos.}

\vfill
\begin{minipage}[c]{0.55\linewidth}
\begin{idea}[Los conceptos básicos]
La \textbf{estadística} recoge, ordena, representa e interpreta datos para sacar conclusiones.\\[3pt]
\textbf{Población}: todos los elementos que queremos estudiar. \ \textbf{Individuo}: cada uno.\\[2pt]
\textbf{Muestra}: la parte de la población que estudiamos cuando no podemos preguntar a todo el mundo. El número de individuos es su \textbf{tamaño}.
\end{idea}
\end{minipage}\hfill
\begin{minipage}[c]{0.42\linewidth}
\begin{exemple}[Ejemplo: \ ¿cuántas horas duermen?]
Preguntamos a 60 de los 600 alumnos de un instituto, elegidos al azar.\\[3pt]
\textbf{Población}: los 600 alumnos. \ \textbf{Individuo}: cada alumno. \ \textbf{Muestra}: los 60 (tamaño 60).
\end{exemple}
\end{minipage}

\vfill
\begin{compte}[Una muestra debe ser representativa]
Se debe parecer a la población. Si solo preguntamos a los de 2.º de Bachillerato la víspera de un examen, las conclusiones no valen para todo el instituto.
\end{compte}

\vfill
\begin{tcbraster}[raster columns=3,raster equal height,raster column skip=8pt]
\begin{tcolorbox}[fmbase,colback=fmA!8,colframe=fmA,coltitle=white,title={Cualitativa}]
Los valores son \textbf{cualidades}, no números.\par\smallskip
\small Color preferido, deporte, tipo de mascota, mes de nacimiento.
\end{tcolorbox}
\begin{tcolorbox}[fmbase,colback=fmB!8,colframe=fmB,coltitle=white,title={Cuantitativa discreta}]
Números que se obtienen \textbf{contando}: valores aislados (0, 1, 2\dots).\par\smallskip
\small Número de hermanos, goles en un partido, libros leídos.
\end{tcolorbox}
\begin{tcolorbox}[fmbase,colback=fmK!8,colframe=fmK,coltitle=white,title={Cuantitativa continua}]
Números que se obtienen \textbf{midiendo}: entre dos valores siempre hay otro.\par\smallskip
\small Estatura, peso, tiempo, temperatura.
\end{tcolorbox}
\end{tcbraster}

\vfill
\begin{truc}[¿Se cuenta o se mide?]
Si se \textbf{cuenta}, es discreta. Si se \textbf{mide} (con regla, reloj, balanza, termómetro\dots), es continua. \ La estatura en cm, $152$; $152{,}5$; $152{,}37$\dots: continua.
\end{truc}

\vfill
\begin{tcbraster}[raster columns=3,raster equal height,raster column skip=8pt]
\begin{tcolorbox}[fmbase,colback=fmdefbg,colframe=fmdef,coltitle=white,title={\faRandom\ \ Aleatorio simple}]
\small Todos tienen la misma probabilidad de salir, como en un sorteo.
\end{tcolorbox}
\begin{tcolorbox}[fmbase,colback=fmdefbg,colframe=fmdef,coltitle=white,title={\faListOl\ \ Sistemático}]
\small El primero al azar y después uno de cada $k$ (uno de cada 10 de la lista).
\end{tcolorbox}
\begin{tcolorbox}[fmbase,colback=fmdefbg,colframe=fmdef,coltitle=white,title={\faLayerGroup\ \ Estratificado}]
\small Por grupos (cursos, edades\dots), de cada uno en proporción a su tamaño.
\end{tcolorbox}
\end{tcbraster}

\vfill
\begin{exemple}[Ejemplo: \ muestra estratificada de 50 en un instituto de 500 alumnos]
\begin{minipage}[c]{0.36\linewidth}
$\dfrac{50}{500}=\dfrac1{10}$: la décima parte de cada curso, elegidos al azar dentro del curso.
\end{minipage}\hfill
\begin{minipage}[c]{0.6\linewidth}\centering\small
\taularod[fmex]{>{\heavy}l*{5}{>{\hspace{5pt}}c<{\hspace{5pt}}}}{
\rowcolor{fmex} & \color{white}1.º & \color{white}2.º & \color{white}3.º & \color{white}4.º & \color{white}Total\\
Alumnos & 150 & 120 & 130 & 100 & 500\\
\rowcolor{fmexbg}Muestra & 15 & 12 & 13 & 10 & 50
}
\end{minipage}
\end{exemple}

\vfill
% ============================================================ 9.2
\clearpage
\seccio{sec2}{Tablas de frecuencias con datos aislados}{Contamos cuántas veces aparece cada valor.}

\vfill
\begin{idea}[Frecuencias]
\textbf{Absoluta} $f_i$: cuántas veces aparece el valor. La suma de todas es $N$, el número de datos.\\[2pt]
\textbf{Relativa} $h_i=\dfrac{f_i}{N}$: la parte del total (fracción, decimal o \%). Todas suman 1 (el 100\,\%).\\[2pt]
\textbf{Acumuladas}: $F_i$, los datos menores o iguales que $x_i$ (sumamos las $f_i$ desde arriba); \ $H_i=\dfrac{F_i}{N}$. La última $F_i$ es $N$.
\end{idea}

\vfill
\begin{exemple}[Ejemplo: \ número de hermanos de 20 alumnos]
\begin{minipage}[c]{0.4\linewidth}
\small 1 \ 2 \ 0 \ 1 \ 1 \ 3 \ 2 \ 1 \ 0 \ 2\\ 1 \ 1 \ 4 \ 0 \ 2 \ 1 \ 3 \ 1 \ 2 \ 1\\[6pt]
\normalsize Recuento con palotes: una raya por dato; la quinta atraviesa las cuatro anteriores.\\[4pt]
9 de los 20 tienen un hermano: \ $h=\frac9{20}=0{,}45$, el \boxed{45\,\%}
\end{minipage}\hfill
\begin{minipage}[c]{0.56\linewidth}\centering\small
\taularod[fmex]{>{\heavy}c<{\hspace{4pt}}>{\hspace{4pt}}l<{\hspace{4pt}}*{3}{>{\hspace{4pt}}c<{\hspace{4pt}}}}{
\rowcolor{fmex}\color{white}$x_i$ & \color{white}\heavy Recuento & \color{white}$f_i$ & \color{white}$h_i$ & \color{white}\%\\
0 & $|||$ & 3 & $0{,}15$ & 15\\
\rowcolor{fmexbg}1 & $\cancel{||||}\ ||||$ & 9 & $0{,}45$ & 45\\
2 & $\cancel{||||}$ & 5 & $0{,}25$ & 25\\
\rowcolor{fmexbg}3 & $||$ & 2 & $0{,}1$ & 10\\
4 & $|$ & 1 & $0{,}05$ & 5\\
\rowcolor{fmexbg}Total & & 20 & 1 & 100
}
\end{minipage}
\end{exemple}

\vfill
\begin{exemple}[Ejemplo: \ notas de un examen de 25 alumnos]
\begin{minipage}[c]{0.56\linewidth}\centering\small
\taularod[fmex]{>{\heavy}c*{6}{>{\hspace{4pt}}c<{\hspace{4pt}}}}{
\rowcolor{fmex}\color{white}$x_i$ & \color{white}$f_i$ & \color{white}$F_i$ & \color{white}$h_i$ & \color{white}$H_i$ & \color{white}\% & \color{white}$\%_A$\\
4 & 2 & 2 & $0{,}08$ & $0{,}08$ & 8 & 8\\
\rowcolor{fmexbg}5 & 5 & 7 & $0{,}2$ & $0{,}28$ & 20 & 28\\
6 & 6 & 13 & $0{,}24$ & $0{,}52$ & 24 & 52\\
\rowcolor{fmexbg}7 & 5 & 18 & $0{,}2$ & $0{,}72$ & 20 & 72\\
8 & 4 & 22 & $0{,}16$ & $0{,}88$ & 16 & 88\\
\rowcolor{fmexbg}9 & 2 & 24 & $0{,}08$ & $0{,}96$ & 8 & 96\\
10 & 1 & 25 & $0{,}04$ & 1 & 4 & 100
}
\end{minipage}\hfill
\begin{minipage}[c]{0.41\linewidth}
Un 6 o menos: \ $F=13$ alumnos, el $\%_A=\boxed{52\,\%}$\\[5pt]
Aprueban (5 o más): \ $25-2=\boxed{23}$\\[5pt]
Más de un 7: \ $25-18=7$ alumnos,\\ el $16+8+4=\boxed{28\,\%}$
\end{minipage}
\end{exemple}

\vfill
\begin{truc}[Comprueba los totales]
Las $f_i$ deben sumar $N$; las $h_i$, 1; los \%, 100; y la última $F_i$ debe ser $N$. Si no cuadra, has perdido algún dato en el recuento.
\end{truc}

\vfill
% ============================================================ 9.3
\clearpage
\seccio{sec3}{Tablas con datos agrupados en intervalos}{Si hay muchos valores distintos, los agrupamos.}

\vfill
\begin{minipage}[c]{0.52\linewidth}
\begin{idea}[Intervalos y marca de clase]
Para variables continuas o con muchos valores hacemos \textbf{intervalos} (clases) de la misma amplitud.\\[3pt]
$[150,\ 155\ro$: entra el 150, pero el 155 va al intervalo siguiente.\\[3pt]
\textbf{Marca de clase}: el punto medio, \ $\dfrac{150+155}{2}=152{,}5$. Representa todos los datos del intervalo.
\end{idea}
\end{minipage}\hfill
\begin{minipage}[c]{0.45\linewidth}
\passosT{Cómo se hacen los intervalos}{\pas{1} \textbf{Rango}: $x_{\max}-x_{\min}$.\\[3pt]
\pas{2} Número de intervalos: entre 5 y 10, cerca de $\sqrt N$.\\[3pt]
\pas{3} \textbf{Amplitud}: rango entre el número de intervalos, redondeado hacia arriba a un número cómodo.\\[3pt]
\pas{4} Empezamos en un valor redondo $\le x_{\min}$.}
\end{minipage}

\vfill
\begin{exemple}[Ejemplo: \ estaturas (cm) de 30 alumnos]
{\small 152 \ 147 \ 158 \ 150 \ 143 \ 161 \ 155 \ 149 \ 151 \ 146 \ 157 \ 153 \ 163 \ 150 \ 148 \ 154 \ 159 \ 141 \ 152 \ 156 \ 145 \ 151 \ 160 \ 149 \ 153 \ 144 \ 158 \ 147 \ 162 \ 150}\\[5pt]
\begin{minipage}[c]{0.36\linewidth}
\pas{1} Rango: \ $163-141=22$ cm\\[3pt]
\pas{2} $\sqrt{30}\approx5{,}5$: \ 5 intervalos\\[3pt]
\pas{3} $22:5=4{,}4\ \to$ amplitud \textbf{5}\\[3pt]
\pas{4} Empezamos en 140
\end{minipage}\hfill
\begin{minipage}[c]{0.62\linewidth}\centering\small
\taularod[fmex]{>{\heavy}c*{5}{>{\hspace{4pt}}c<{\hspace{4pt}}}}{
\rowcolor{fmex}\color{white}Estatura & \color{white}Marca & \color{white}$f_i$ & \color{white}$F_i$ & \color{white}$h_i$ & \color{white}\%\\
$[140,145\ro$ & 142,5 & 3 & 3 & $0{,}1$ & 10\\
\rowcolor{fmexbg}$[145,150\ro$ & 147,5 & 7 & 10 & $0{,}233$ & 23,3\\
$[150,155\ro$ & 152,5 & 10 & 20 & $0{,}333$ & 33,3\\
\rowcolor{fmexbg}$[155,160\ro$ & 157,5 & 6 & 26 & $0{,}2$ & 20\\
$[160,165\ro$ & 162,5 & 4 & 30 & $0{,}133$ & 13,3
}
\end{minipage}
\end{exemple}

\vfill
\begin{minipage}[t]{0.5\linewidth}
\begin{compte}[¿Dónde va el 150?]
Los tres alumnos que miden exactamente 150 cm se cuentan en $[150,155\ro$, no en $[145,150\ro$.
\end{compte}
\end{minipage}\hfill
\begin{minipage}[t]{0.47\linewidth}
\begin{truc}[Agrupar pierde información]
Dentro de un intervalo ya no sabemos los datos exactos: los cálculos con las marcas de clase serán \textbf{aproximados}.
\end{truc}
\end{minipage}

\vfill
\begin{exemple}[Ejemplo: \ dinero semanal de 40 alumnos, en intervalos de 5 €]
\centering\small
\taularod[fmex]{>{\heavy}c*{5}{>{\hspace{6pt}}c<{\hspace{6pt}}}}{
\rowcolor{fmex}\color{white}Dinero (€) & \color{white}$[0,5\ro$ & \color{white}$[5,10\ro$ & \color{white}$[10,15\ro$ & \color{white}$[15,20\ro$ & \color{white}$[20,25\ro$\\
Marca & 2,5 & 7,5 & 12,5 & 17,5 & 22,5\\
\rowcolor{fmexbg}$f_i$ & 6 & 10 & 12 & 8 & 4\\
$F_i$ & 6 & 16 & 28 & 36 & 40
}
\end{exemple}

\vfill
% ============================================================ 9.4
\clearpage
\seccio{sec4}{Gráficos estadísticos}{De un vistazo, cómo se reparten los datos.}

\vfill
\begin{tcbraster}[raster columns=2,raster equal height,raster column skip=10pt]
\begin{tcolorbox}[fmbase,colback=fmdefbg,colframe=fmdef,coltitle=white,title={Diagrama de barras y polígono},before upper=\centering]
\begin{tikzpicture}[x=0.85cm,y=0.3cm]
  \draw[-{Latex[length=2mm]},fmmuted,line width=0.8pt] (-0.6,0) -- (4.7,0);
  \draw[-{Latex[length=2mm]},fmmuted,line width=0.8pt] (-0.6,0) -- (-0.6,10.5) node[above,font=\scriptsize] {$f_i$};
  \foreach \k in {2,4,...,10} {\draw[fmmuted!30] (-0.6,\k) -- (4.5,\k); \node[left,font=\scriptsize] at (-0.6,\k) {\k};}
  \foreach \x/\f in {0/3,1/9,2/5,3/2,4/1} {\barra{\x}{\f}{fmA} \node[below,font=\scriptsize] at (\x,0) {\x};}
  \draw[fmR,line width=1.4pt] (0,3) -- (1,9) -- (2,5) -- (3,2) -- (4,1); \foreach \x/\f in {0/3,1/9,2/5,3/2,4/1} \fill[fmR] (\x,\f) circle (2.5pt);
\end{tikzpicture}\par\smallskip
\small Hermanos de 20 alumnos. Para variables \textbf{cualitativas} y \textbf{discretas}: barras separadas.
\end{tcolorbox}
\begin{tcolorbox}[fmbase,colback=fmdefbg,colframe=fmdef,coltitle=white,title={Histograma},before upper=\centering]
\begin{tikzpicture}[x=0.15cm,y=0.27cm]
  \draw[-{Latex[length=2mm]},fmmuted,line width=0.8pt] (140,0) -- (167,0) node[right,font=\scriptsize] {cm};
  \draw[-{Latex[length=2mm]},fmmuted,line width=0.8pt] (140,0) -- (140,11.5) node[above,font=\scriptsize] {$f_i$};
  \foreach \k in {2,4,...,10} {\draw[fmmuted!30] (140,\k) -- (165,\k); \node[left,font=\scriptsize] at (140,\k) {\k};}
  \foreach \a/\f in {140/3,145/7,150/10,155/6,160/4} {\filldraw[fill=fmK!30,draw=fmK!80!black,line width=0.9pt] (\a,0) rectangle (\a+5,\f);}
  \foreach \a in {140,150,160} \node[below,font=\scriptsize] at (\a,0) {\a};
  \draw[fmR,line width=1.4pt] (142.5,3) -- (147.5,7) -- (152.5,10) -- (157.5,6) -- (162.5,4);
\end{tikzpicture}\par\smallskip
\small Estaturas de 30 alumnos. Datos en \textbf{intervalos}: barras juntas; el polígono, por las marcas de clase.
\end{tcolorbox}
\end{tcbraster}

\vfill
\begin{compte}[Un gráfico bien hecho]
Pon siempre el \textbf{título} y el nombre de cada \textbf{eje}, y empieza el eje vertical en \textbf{0}: si no, las diferencias entre barras parecen mayores de lo que son.
\end{compte}

\vfill
\begin{minipage}[c]{0.4\linewidth}\centering
\begin{tikzpicture}[scale=1.25]
  \filldraw[fill=fmA!45,draw=white,line width=1.2pt] (0,0) -- (0:1.4) arc (0:150:1.4) -- cycle;
  \filldraw[fill=fmB!45,draw=white,line width=1.2pt] (0,0) -- (150:1.4) arc (150:240:1.4) -- cycle;
  \filldraw[fill=fmK!45,draw=white,line width=1.2pt] (0,0) -- (240:1.4) arc (240:315:1.4) -- cycle;
  \filldraw[fill=fmV!45,draw=white,line width=1.2pt] (0,0) -- (315:1.4) arc (315:360:1.4) -- cycle;
  \node[font=\scriptsize] at (75:0.85) {andando}; \node[font=\scriptsize] at (195:0.85) {autobús};
  \node[font=\scriptsize] at (277:0.85) {coche}; \node[font=\scriptsize] at (337:1.05) {bici};
\end{tikzpicture}
\end{minipage}\hfill
\begin{minipage}[c]{0.57\linewidth}
\begin{idea}[Diagrama de sectores]
El círculo es el total; cada sector, proporcional a su frecuencia:\\[3pt]
{\centering$\text{ángulo}=\dfrac{f_i}{N}\cdot360^\circ=h_i\cdot360^\circ$\par}\smallskip
Sobre todo para variables \textbf{cualitativas}: qué parte del total es cada valor.
\end{idea}
\end{minipage}

\vfill
\begin{exemple}[Ejemplo: \ cómo vienen al instituto los 24 alumnos de una clase]
\begin{minipage}[c]{0.55\linewidth}\centering\small
\taularod[fmex]{>{\heavy}l*{4}{>{\hspace{5pt}}c<{\hspace{5pt}}}}{
\rowcolor{fmex} & \color{white}Andando & \color{white}Autobús & \color{white}Coche & \color{white}Bici\\
$f_i$ & 10 & 6 & 5 & 3\\
\rowcolor{fmexbg}Ángulo & $150^\circ$ & $90^\circ$ & $75^\circ$ & $45^\circ$
}
\end{minipage}\hfill
\begin{minipage}[c]{0.42\linewidth}
\textbf{Truco}: $360:24=15^\circ$ por alumno.\\ Andando: $10\cdot15^\circ=150^\circ$.\\[4pt]
\textbf{Pictograma}: un dibujo para cada cantidad (\faBook\ = 10 libros); medio dibujo, 5.
\end{minipage}
\end{exemple}

\vfill
% ============================================================ 9.5
\clearpage
\seccio{sec5}{Parámetros de centralización}{Un solo valor «del medio» que resume todos los datos.}

\vfill
\begin{tcbraster}[raster columns=3,raster equal height,raster column skip=8pt]
\begin{tcolorbox}[fmbase,colback=fmA!8,colframe=fmA,coltitle=white,title={Moda $Mo$}]
El valor que \textbf{más se repite}. El único para variables cualitativas.
\end{tcolorbox}
\begin{tcolorbox}[fmbase,colback=fmB!8,colframe=fmB,coltitle=white,title={Media $\bar x$}]
La suma de los datos entre $N$: el \textbf{reparto a partes iguales}.
\end{tcolorbox}
\begin{tcolorbox}[fmbase,colback=fmK!8,colframe=fmK,coltitle=white,title={Mediana $Me$}]
El valor del \textbf{centro} con los datos ordenados.
\end{tcolorbox}
\end{tcbraster}

\vfill
\begin{exemple}[Ejemplo: \ minutos que tardan 7 alumnos en llegar al instituto: \ 20, 40, 15, 46, 20, 30, 25]
\begin{minipage}[c]{0.6\linewidth}
$Mo=\boxed{20}$ min \ \nota{aparece dos veces}\\[4pt]
$\bar x=\dfrac{20+40+15+46+20+30+25}{7}=\dfrac{196}{7}=\boxed{28}$ min\\[4pt]
Ordenados: \ $15,\ 20,\ 20,\ \boxed{25},\ 30,\ 40,\ 46$: \ $Me=\boxed{25}$ min
\end{minipage}\hfill
\begin{minipage}[c]{0.36\linewidth}
Con $N$ par: \ $3,\ 5,\ 6,\ \boxed{7,\ 8},\ 8,\ 9,\ 10$\\[3pt]
$Me=\dfrac{7+8}{2}=7{,}5$
\end{minipage}
\end{exemple}

\vfill
\begin{compte}[La media y los datos extremos]
Los datos extremos mueven mucho la media, pero no la mediana: si los hay, la \textbf{mediana} representa mejor el conjunto. Y antes de buscar la mediana, ¡\textbf{ordena}!
\end{compte}

\vfill
\begin{exemple}[Ejemplo: \ con una tabla (hermanos de 20 alumnos)]
\begin{minipage}[c]{0.42\linewidth}\centering\small
\taularod[fmex]{>{\heavy}c*{3}{>{\hspace{5pt}}c<{\hspace{5pt}}}}{
\rowcolor{fmex}\color{white}$x_i$ & \color{white}$f_i$ & \color{white}$F_i$ & \color{white}$x_i\cdot f_i$\\
0 & 3 & 3 & 0\\
\rowcolor{fmexbg}1 & 9 & 12 & 9\\
2 & 5 & 17 & 10\\
\rowcolor{fmexbg}3 & 2 & 19 & 6\\
4 & 1 & 20 & 4\\
\rowcolor{fmexbg}Total & 20 & & 29
}
\end{minipage}\hfill
\begin{minipage}[c]{0.55\linewidth}
$\bar x=\dfrac{\sum x_i\cdot f_i}{N}=\dfrac{29}{20}=\boxed{1{,}45}$ hermanos\\[4pt]
$Me$: centrales el 10.º y el 11.º; la primera $F_i\ge10$ es 12, del 1: \ $Me=\boxed{1}$\\[4pt]
$Mo=\boxed{1}$ \ \nota{$f=9$, la mayor}\\[3pt]
{\small Nadie tiene 1,45 hermanos: es el reparto a partes iguales.}
\end{minipage}
\end{exemple}

\vfill
\begin{exemple}[Ejemplo: \ datos agrupados (dinero de 40 alumnos, 9.3), con las marcas de clase]
$\bar x\approx\dfrac{2{,}5\cdot6+7{,}5\cdot10+12{,}5\cdot12+17{,}5\cdot8+22{,}5\cdot4}{40}=\dfrac{470}{40}=\boxed{11{,}75\text{ €}}$\\[4pt]
Intervalo \textbf{modal}: $[10,15\ro$ \nota{$f=12$} \qquad Intervalo \textbf{mediano}: $\frac N2=20$, la primera $F_i\ge20$ es 28: \ $[10,15\ro$
\end{exemple}

\vfill
\begin{truc}[Marta: 6, 7 y 5. ¿Qué nota necesita para un 6,5 de media en 4 exámenes?]
$\text{suma}=\bar x\cdot N$: \ $6{,}5\cdot4=26$; \ ya tiene 18 \ $\to$ \ $\boxed{8}$. \ Para juntar grupos, sumamos las sumas (no la media de las medias).
\end{truc}

\vfill
% ============================================================ 9.6
\clearpage
\seccio{sec6}{Parámetros de dispersión}{Cómo de dispersos están los datos alrededor de la media.}

\vfill
\begin{idea}[Misma media, datos muy diferentes]
Puntos en cinco partidos: \ \textbf{Anna} $8,\,9,\,10,\,11,\,12$ \ y \ \textbf{Bea} $2,\,6,\,10,\,14,\,18$. Las dos tienen una media de 10, pero Anna es mucho más regular.
\end{idea}

\vfill
\begin{minipage}[c]{0.54\linewidth}\small
\textbf{Rango}: \ $x_{\max}-x_{\min}$. Fácil, pero solo depende de dos datos.\\[4pt]
\textbf{Desviación media}: \ $DM=\dfrac{\sum|x_i-\bar x|\cdot f_i}{N}$\\[4pt]
\textbf{Varianza}: \ $\sigma^2=\dfrac{\sum(x_i-\bar x)^2\cdot f_i}{N}$\\[4pt]
\textbf{Desviación típica}: \ $\sigma=\sqrt{\sigma^2}$, con las mismas unidades que los datos: la más usada.\\[4pt]
Cuanto mayores, más dispersos y \textbf{menos representativa} es la media.
\end{minipage}\hfill
\begin{minipage}[c]{0.43\linewidth}\centering
\begin{tikzpicture}[x=0.3cm]
  \draw[fmmuted] (1,1) -- (19,1); \draw[fmmuted] (1,0) -- (19,0);
  \foreach \x in {8,9,10,11,12} \fill[fmA] (\x,1) circle (4pt);
  \foreach \x in {2,6,10,14,18} \fill[fmB] (\x,0) circle (4pt);
  \draw[fmR,dashed,line width=1pt] (10,-0.5) -- (10,1.5) node[above,font=\scriptsize] {$\bar x=10$};
  \node[left,font=\scriptsize\heavy,text=fmA] at (1,1) {Anna}; \node[left,font=\scriptsize\heavy,text=fmB] at (1,0) {Bea};
\end{tikzpicture}
\end{minipage}

\vfill
\begin{exemple}[Ejemplo: \ Anna y Bea]
\begin{minipage}[c]{0.4\linewidth}\centering\small
\taularod[fmex]{>{\heavy}c*{3}{>{\hspace{4pt}}c<{\hspace{4pt}}}}{
\rowcolor{fmex}\color{white}$x_i$ & \color{white}$x_i-\bar x$ & \color{white}$|x_i-\bar x|$ & \color{white}$(x_i-\bar x)^2$\\
8 & $-2$ & 2 & 4\\
\rowcolor{fmexbg}9 & $-1$ & 1 & 1\\
10 & 0 & 0 & 0\\
\rowcolor{fmexbg}11 & 1 & 1 & 1\\
12 & 2 & 2 & 4\\
\rowcolor{fmexbg}$\sum$ & & 6 & 10
}
\end{minipage}\hfill
\begin{minipage}[c]{0.57\linewidth}
\textbf{Anna}: rango 4; \ $DM=\frac65=1{,}2$; \ $\sigma^2=\frac{10}5=2$; \ $\sigma=\sqrt2\approx\boxed{1{,}41}$\\[5pt]
\textbf{Bea}: desviaciones $-8,-4,0,4,8$: \ $\sum|\ |=24$, \ $\sum(\ )^2=160$\\
rango 16; \ $DM=4{,}8$; \ $\sigma^2=32$; \ $\sigma=\sqrt{32}\approx\boxed{5{,}66}$\\[5pt]
Todos dicen lo mismo: los de Bea están mucho más dispersos.
\end{minipage}
\end{exemple}

\vfill
\begin{tcolorbox}[enhanced,colback=fmyellow!18,colframe=fmgold,arc=4mm,boxrule=1.6pt,before upper=\centering,top=2mm,bottom=2mm]
{\Large$\sigma^2=\dfrac{\sum x_i^2\cdot f_i}{N}-\bar x^{\,2}$}\par\smallskip
{\small la media de los cuadrados menos el cuadrado de la media: más cómoda si $\bar x$ no es entera}
\end{tcolorbox}

\vfill
\begin{exemple}[Ejemplo: \ hermanos de 20 alumnos]
\begin{minipage}[c]{0.4\linewidth}\centering\small
\taularod[fmex]{>{\heavy}c*{3}{>{\hspace{5pt}}c<{\hspace{5pt}}}}{
\rowcolor{fmex}\color{white}$x_i$ & \color{white}$f_i$ & \color{white}$x_if_i$ & \color{white}$x_i^2f_i$\\
0 & 3 & 0 & 0\\
\rowcolor{fmexbg}1 & 9 & 9 & 9\\
2 & 5 & 10 & 20\\
\rowcolor{fmexbg}3 & 2 & 6 & 18\\
4 & 1 & 4 & 16\\
\rowcolor{fmexbg}$\sum$ & 20 & 29 & 63
}
\end{minipage}\hfill
\begin{minipage}[c]{0.57\linewidth}
$\bar x=\dfrac{29}{20}=1{,}45$\\[4pt]
$\sigma^2=\dfrac{63}{20}-1{,}45^2=3{,}15-2{,}1025=1{,}0475$\\[4pt]
$\sigma=\sqrt{1{,}0475}\approx\boxed{1{,}02}$ hermanos\\[3pt]
\nota{una varianza negativa significa que te has equivocado}
\end{minipage}
\end{exemple}

\vfill
\begin{exemple}[Ejemplo: \ datos agrupados (dinero de 40 alumnos, 9.3)]
Con las marcas de clase: \ $\sum x_i^2f_i=37{,}5+562{,}5+1875+2450+2025=6950$; \ $\bar x=11{,}75$\\[4pt]
$\sigma^2=\dfrac{6950}{40}-11{,}75^2=173{,}75-138{,}06=35{,}69$ \qquad $\sigma\approx\boxed{5{,}97\text{ €}}$ \qquad $CV=\dfrac{5{,}97}{11{,}75}\approx\boxed{51\,\%}$ \ \nota{muy dispersos}
\end{exemple}

\vfill
\begin{minipage}[c]{0.5\linewidth}
\begin{idea}[Coeficiente de variación]
{\centering\Large$CV=\dfrac{\sigma}{\bar x}$\par}\smallskip
Sin unidades (a menudo en \%): compara la dispersión de grupos con medias o unidades muy diferentes.
\end{idea}
\end{minipage}\hfill
\begin{minipage}[c]{0.47\linewidth}
\begin{exemple}[Ejemplo: \ bebés y adultos]
Bebés: $\bar x=3{,}2$, $\sigma=0{,}4$ kg: \ $CV=\boxed{12{,}5\,\%}$\\[3pt]
Adults: $\bar x=70$, $\sigma=8$ kg: \ $CV\approx\boxed{11{,}4\,\%}$\\[3pt]
En proporción, los bebés tienen \textbf{más} dispersión.
\end{exemple}
\end{minipage}

\vfill
\begin{tcbraster}[raster columns=2,raster equal height,raster column skip=10pt]
\begin{exemple}[Ejemplo: \ dos marcas de bombillas]
A: $\bar x=1000$ h, $\sigma=50$ h \quad B: $\bar x=1000$ h, $\sigma=200$ h\\[3pt]
Duran lo mismo de media, pero las de \textbf{A} son mucho más regulares: es más fácil saber cuánto durarán.
\end{exemple}
\begin{truc}[Con la calculadora]
En modo estadístico, nuestra $\sigma$ es la que se llama $\sigma_x$ (o $\sigma_n$). La $s_x$ divide entre $N-1$: se usa en cursos posteriores.
\end{truc}
\end{tcbraster}

\vfill
% ============================================================ 9.7
\clearpage
\seccio{sec7}{Posición: cuartiles y percentiles}{Qué lugar ocupa un dato dentro de los datos ordenados.}

\vfill
\begin{minipage}[c]{0.55\linewidth}
\begin{idea}[Cuartiles y percentiles]
\textbf{Cuartiles}: dividen los datos ordenados en cuatro partes iguales. $Q_1$ deja por debajo el 25\,\%; \ $Q_2=Me$, el 50\,\%; \ $Q_3$, el 75\,\%.\\[3pt]
\textbf{Percentil} $P_k$: deja por debajo el $k$\,\%. \ $Q_1=P_{25}$, $Me=P_{50}$, $Q_3=P_{75}$.\\[3pt]
\textbf{Rango intercuartílico}: $RI=Q_3-Q_1$, donde está el 50\,\% central. No lo mueven los datos extremos.
\end{idea}
\end{minipage}\hfill
\begin{minipage}[c]{0.42\linewidth}
\passosT{Cómo se calculan}{\pas{1} \textbf{Ordenamos} los $N$ datos.\\[3pt]
\pas{2} Posición: \ $\dfrac{k\cdot N}{100}$ \ ($Q_1$: $\frac N4$; \ $Q_3$: $\frac{3N}4$).\\[3pt]
\pas{3} \textbf{No entera}: el dato de la posición siguiente. \textbf{Entera}: la media de ese dato y el siguiente.}
\end{minipage}

\vfill
\begin{exemple}[Ejemplo: \ notas de 12 alumnos, ya ordenadas]
\centering\small
\taularod[fmex]{>{\heavy}l*{12}{>{\hspace{3pt}}c<{\hspace{3pt}}}}{
\rowcolor{fmex}\color{white}Posición & \color{white}1 & \color{white}2 & \color{white}3 & \color{white}4 & \color{white}5 & \color{white}6 & \color{white}7 & \color{white}8 & \color{white}9 & \color{white}10 & \color{white}11 & \color{white}12\\
Dato & 3 & 4 & \cellcolor{fmA!20}5 & \cellcolor{fmA!20}5 & 6 & \cellcolor{fmK!20}6 & \cellcolor{fmK!20}7 & 7 & \cellcolor{fmB!20}8 & \cellcolor{fmB!20}8 & 9 & 10
}\par\medskip\raggedright\normalsize
$Q_1$: $\frac{12}4=3$, entera: \ $\frac{5+5}2=\boxed{5}$ \qquad $Me$: $\frac{12}2=6$: \ $\frac{6+7}2=\boxed{6{,}5}$ \qquad $Q_3$: $\frac{36}4=9$: \ $\frac{8+8}2=\boxed{8}$ \qquad $RI=8-5=3$\\[4pt]
$P_{90}$: \ $\frac{90\cdot12}{100}=10{,}8$, no entera: el dato 11.º, \ $P_{90}=\boxed{9}$
\end{exemple}

\vfill
\begin{exemple}[Ejemplo: \ con una tabla (notas de 25 alumnos, 9.2): \ la primera $F_i$ que llega a la posición]
\begin{minipage}[c]{0.4\linewidth}\centering\small
\taularod[fmex]{>{\heavy}c*{7}{>{\hspace{2pt}}c<{\hspace{2pt}}}}{
\rowcolor{fmex}\color{white}$x_i$ & \color{white}4 & \color{white}5 & \color{white}6 & \color{white}7 & \color{white}8 & \color{white}9 & \color{white}10\\
$F_i$ & 2 & 7 & 13 & 18 & 22 & 24 & 25
}
\end{minipage}\hfill
\begin{minipage}[c]{0.57\linewidth}
$Q_1$: $\frac{25}4=6{,}25\to F=7$: \ $Q_1=\boxed{5}$\\[3pt] $Me$: $12{,}5\to F=13$: \ $Me=\boxed{6}$\\[3pt]
$Q_3$: $18{,}75\to F=22$: \ $Q_3=\boxed{8}$\\[3pt] $P_{90}$: $22{,}5\to F=24$: \ $P_{90}=\boxed{9}$\\[3pt]
\nota{si una $F_i$ coincide exactamente con la posición, media de ese valor y el siguiente}
\end{minipage}
\end{exemple}

\vfill
\begin{truc}[Percentiles en la vida real]
Un bebé en el \textbf{percentil 80} de peso pesa más que el 80\,\% de los bebés de su edad (y menos que el 20\,\%). Con datos agrupados, el cuartil o percentil está en el primer intervalo con $F_i\ge\frac{kN}{100}$.
\end{truc}

\vfill
% ============================================================ 9.8
\clearpage
\seccio{sec8}{Diagrama de cajas}{Cinco números que resumen toda la distribución.}

\vfill
\begin{minipage}[c]{0.5\linewidth}
\begin{idea}[El diagrama de cajas (caja y bigotes)]
\pas{1} Una escala horizontal con todos los datos.\\[3pt]
\pas{2} Una \textbf{caja} de $Q_1$ a $Q_3$, con la \textbf{mediana} marcada.\\[3pt]
\pas{3} Los \textbf{bigotes}: de la caja hasta el mínimo y hasta el máximo.
\end{idea}
\end{minipage}\hfill
\begin{minipage}[c]{0.47\linewidth}\centering
\begin{tikzpicture}[x=0.75cm]
  \draw[fmmuted,line width=0.8pt] (3.5,0) -- (10.5,0); \foreach \k in {4,...,10} {\draw[fmmuted] (\k,-0.08) -- (\k,0.08); \node[below,font=\scriptsize] at (\k,-0.1) {\k};}
  \caixa{1}{4}{5}{6}{8}{10}{fmblue}
  \foreach \x/\t in {4/mín,5/$Q_1$,6/$Me$,8/$Q_3$,10/máx} \node[above,font=\scriptsize] at (\x,1.5) {\t};
  \foreach \a/\b in {4/5,5/6,6/8,8/10} \node[font=\scriptsize,text=fmR] at ({(\a+\b)/2},2.05) {25\,\%};
\end{tikzpicture}\\
{\comic\small\color{fmmuted}notas de 25 alumnos (9.7)}
\end{minipage}

\vfill
\begin{tcbraster}[raster columns=2,raster equal height,raster column skip=10pt]
\begin{tcolorbox}[fmbase,colback=fmdefbg,colframe=fmdef,coltitle=white,title={Cómo se interpreta}]
Cada parte (bigote, media caja, media caja, bigote) tiene el \textbf{25\,\%} de los datos. Una parte \textbf{corta}: datos concentrados; \textbf{larga}: dispersos.\par\smallskip
La caja mide el $RI$; el total, el rango.
\end{tcolorbox}
\begin{tcolorbox}[fmbase,colback=fmdefbg,colframe=fmdef,coltitle=white,title={Simetría}]
Mediana en el centro de la caja y bigotes iguales: \textbf{simétrica}. Si no, \textbf{asimétrica}, hacia el lado de la parte más larga.\par\smallskip
En el ejemplo, de 5 a 6 muy juntos; de 6 a 10, dispersos.
\end{tcolorbox}
\end{tcbraster}

\vfill
\begin{exemple}[Ejemplo: \ dos clases con la misma media (6,75) y la misma mediana (7)]
\begin{minipage}[c]{0.5\linewidth}\centering
\begin{tikzpicture}[x=0.6cm,y=0.9cm]
  \draw[fmmuted,line width=0.8pt] (1.5,0) -- (10.5,0); \foreach \k in {2,...,10} {\draw[fmmuted] (\k,-0.06) -- (\k,0.06); \node[below,font=\scriptsize] at (\k,-0.1) {\k};}
  \caixa{2}{4}{5.5}{7}{8}{9}{fmA} \node[left,font=\small\heavy,text=fmA] at (1.5,2) {A};
  \caixa{0.9}{2}{4.5}{7}{9.5}{10}{fmB} \node[left,font=\small\heavy,text=fmB] at (1.5,0.9) {B};
\end{tikzpicture}
\end{minipage}\hfill
\begin{minipage}[c]{0.47\linewidth}
$RI_A=8-5{,}5=2{,}5$ \qquad $RI_B=9{,}5-4{,}5=5$\\[4pt]
La caja de B es el doble de larga: las notas de la clase \textbf{A} son mucho más \textbf{homogéneas}; en la B las hay muy bajas y muy altas.\\[3pt]
\nota{para comparar, siempre con la misma escala}
\end{minipage}
\end{exemple}

\vfill
\begin{exemple}[Para saber más: \ valores atípicos]
\begin{minipage}[c]{0.55\linewidth}
Es \textbf{atípico} un dato menor que $Q_1-1{,}5\cdot RI$ o mayor que $Q_3+1{,}5\cdot RI$.\\[4pt]
$12,\,14,\,15,\,15,\,16,\,17,\,18,\,19,\,35$: \ $Q_1=15$, $Q_3=18$, $RI=3$\\[3pt]
Límites: $10{,}5$ y $22{,}5$: \ el \boxed{35} es atípico.
\end{minipage}\hfill
\begin{minipage}[c]{0.42\linewidth}\centering
\begin{tikzpicture}[x=0.17cm]
  \draw[fmmuted,line width=0.8pt] (10,0) -- (37,0); \foreach \k in {10,15,...,35} {\draw[fmmuted] (\k,-0.08) -- (\k,0.08); \node[below,font=\scriptsize] at (\k,-0.1) {\k};}
  \caixa{0.8}{12}{15}{16}{18}{19}{fmblue} \fill[fmR] (35,0.8) circle (3pt) node[above,font=\scriptsize,text=fmR] {atípico};
\end{tikzpicture}
\end{minipage}
\end{exemple}

\vfill
% ============================================================ resum
\resumtitol
\begin{tcbraster}[raster columns=2,raster equal height=rows,raster column skip=14pt,raster row skip=30pt,raster before skip=16pt]
\rtile{9.1}{Población y muestra}{Muestra representativa, al azar. Variables: cualitativa, discreta (se cuenta) o continua (se mide).}
\rtile{9.2}{Frecuencias}{$f_i$, \ $h_i=\frac{f_i}{N}$, \ $F_i$ acumulada. La última $F_i$ es $N$.}
\rtile{9.3}{Intervalos}{$[a,b\ro$. Marca de clase: el punto medio.\\[1ex] Amplitud $\approx$ rango $:\sqrt N$.}
\rtile{9.4}{Gráficos}{Barras, histograma, polígono.\\[1ex] Sectores: $h_i\cdot360^\circ$.}
\rtile{9.5}{Centralización}{$Mo$, \ $\bar x=\frac{\sum x_if_i}{N}$, \ $Me$: el del centro.}
\rtile{9.6}{Dispersión}{Rango, $DM$, \ $\sigma^2=\frac{\sum x_i^2f_i}{N}-\bar x^2$, \ $\sigma$, \ $CV=\frac{\sigma}{\bar x}$.}
\rtile{9.7}{Posición}{$Q_1$, $Me$, $Q_3$, $P_k$: posición $\frac{kN}{100}$. \ $RI=Q_3-Q_1$.}
\rtile{9.8}{Cajas}{Mín, $Q_1$, $Me$, $Q_3$, máx. Cada parte, el 25\,\%.}
\end{tcbraster}
\end{document}
